\section{Introduction}\label{sec:intro}

\subsection{The problem}

Erd\H{o}s and Graham \cite{ErGr80} asked (Erd\H{o}s problem 396) for which
$k$ there exist infinitely many $n$ such that
$\prod_{i=0}^{k}(n-i)\mid\binom{2n}{n}$. For the single shifted factor,
Pomerance \cite{Pom15} showed that $(n-k)\mid\binom{2n}{n}$ has
$\gg x/\log^2 x$ witnesses up to $x$ for every fixed $k$, and studied the
\emph{governor set}
\[
\D \;:=\; \Big\{m\ge1:\ m\,\Big|\,\binom{2m}{m}\Big\},
\]
whose membership is, by Kummer's theorem, a condition on the base-$p$
digits of $m$ at the primes $p\mid m$
(Proposition~\ref{prop:kummer}). Ford and Konyagin \cite{FK21} proved that
$\D$ has asymptotic density
\[
c_1 \;=\; 0.11424\ldots,
\]
and developed the digit-counting and exponential-sum machinery on which the
present program rests.

We study the first genuinely \emph{binary} case,
\[
\W \;:=\; \Big\{n\ge2:\ n(n-1)\,\Big|\,\binom{2n}{n}\Big\}
\;=\;\{2,\,210,\,460,\,\ldots\},
\]
where two consecutive integers must divide the same central binomial
coefficient. To our knowledge even the infinitude of $\W$ is not in the
literature.

\subsection{Results and targets}

The program's targets, in decreasing strength:

\begin{theorem}[Target 1$''$; rung i]\label{thm:main}
\statusO\quad
At a set of scales $x$ of logarithmic density $1$,
\[
\frac{1}{\log x}\sum_{n\le x}\frac{\mathbf 1_{\W}(n)}{n}
\;=\; c_1^2 + O\big((\log x)^{-c}\big)
\]
for some $c>0$. In particular $\W$ is infinite and has logarithmic density
$c_1^2$.
\end{theorem}

\begin{theorem}[Target 1$'$; rung i-minus]\label{thm:lower}
\statusO\quad
$\W$ has positive lower logarithmic density.
\end{theorem}

\begin{theorem}[Conditional form; rung ii]\label{thm:conditional}
\textbf{[REDUCTION IN PLACE; see \S\ref{sec:remaining} for the precise
hypothesis list]}\quad
Theorem~\ref{thm:main} holds conditionally on Estimate~D$^\dagger$
(Estimate~\ref{est:Ddagger}) together with Lemmas B and C
(\S\ref{sec:architecture}).
\end{theorem}

The plan of proof is a conditional-independence cascade. Lemma~0
(\S\ref{sec:sandwich}) replaces $\W$, \emph{exactly}, by the correlation
set $\D\cap(\D^{(\pm)}+1)$ of two fixed digit-defined sets, so that
Theorem~\ref{thm:main} becomes a binary independence statement. That
statement then factors through: Lemma~A (small primes are ignorable),
Lemma~B (independence of the \emph{anatomies} --- the large-prime
factorization shapes --- of $n$ and $n-1$), Lemma~C (equidistribution of the
$n$-side digit data against multiplicative weights on $n-1$), and Lemma~D
(cross-side digit independence), the hard core. \S\ref{sec:architecture}
states all four precisely.

The hard core has been reduced (\S\S\ref{sec:dls}--\ref{sec:anatomy}) to:
a proved deep large sieve for Fermat-quotient characters covering the deep
harmonics; a two-frequency minor/major-arc dichotomy in which the minor
arcs are proved (Theorem~\ref{thm:minor}) and the major arcs are counted
subject to two average-equidistribution steps
(Hypotheses~\ref{hyp:E1},~\ref{hyp:E2}); and an exact elementary unwinding
of the anatomy layer into classical Type-I/II and primes-in-progressions
sums (\S\ref{sec:anatomy}). No structural obstruction is currently
identified; the residue is quantitative bookkeeping, listed exhaustively in
\S\ref{sec:remaining}.

\subsection{Status conventions}

Every numbered statement carries a tag: \statusP{} (or \statusPM{} when a
numerical verification at scale is also on record), \statusC, \statusKA{}
(a known proof adapts with low risk), \statusN{} (new but
standard-toolbox), or \statusO. The tags are the paper's contract with
itself: Theorem~\ref{thm:conditional} is claimed only modulo the statements
tagged open, each of which appears verbatim in the ledger of
\S\ref{sec:remaining}.

\subsection{Notation}

$p,q,\ell$ denote primes. $\fr{y}$ is the fractional part, $\nm{y}$ the
distance from $y$ to $\Z$, $\e{y}=e^{2\pi iy}$. $P^+(n)$ is the largest
prime factor of $n$. For $x$ the basic scale we write $\Llog=\log x$.
Throughout, $n$ ranges near $x$ and the \emph{band} refers to prime factors
$p=x^{u}$ with $u\in(1/3,1/2)$; band primes carry the single binding digit
condition (\S\ref{sec:architecture}). $A\ll B$ and $A=O(B)$ are Vinogradov
and Landau notation; subscripts indicate parameter dependence.
